\section{Endpoint geometry and conditional uniqueness}
\label{sec:endpoints}
The physical caps give useful boundary data for the shape metric.
One endpoint coalesces the critical points; the other sends eight
vacua to infinity in the original coordinates. Both must be treated
in a fixed source normalization.

\subsection{The homogeneous endpoint}
Write
\[
f_\epsilon(Z,W)=\frac{Z^4+W^4+Z^2W^2}{16}
                 +\frac{\epsilon}{8}(Z^2+W^2),
\qquad \epsilon=c^{-1/2}.
\]
This follows from $Z=c^{1/4}A$, $W=c^{1/4}B$. The leading cubic
gradient has no zero on the unit sphere: its equations are
$Z(2Z^2+W^2)=W(2W^2+Z^2)=0$, whose only common solution is zero.
Uniformly for bounded $\epsilon$, with $R^2=|Z|^2+|W|^2$,
\[
|\nabla f_\epsilon|\ge aR^3-bR,\qquad
|\nabla^k f_\epsilon|\le C_k(1+R^{4-k}),\quad k\ge2.
\]
It follows that for every $\delta>0$ and relevant $k$,
$\delta|\nabla f_\epsilon|^k-|\nabla^k f_\epsilon|\to+\infty$.
Euclidean space has bounded geometry, so these estimates establish
the strong ellipticity hypotheses of the non-Morse Hodge
comparison theorem \cite[Theorem 2.34]{LiWen}.

\begin{proposition}\label{end:continuous}
The harmonic spaces of $f_\epsilon$, including $\epsilon=0$, have
dimension nine. Their projections and the Gram matrices in fixed
compact cohomology representatives vary continuously at zero.
\end{proposition}
\begin{proof}
The two leading cubics form a regular sequence. Their quotient has
Hilbert series $(1+z+z^2)^2$ and dimension nine. Lower-degree terms
preserve this dimension, and Hodge comparison identifies it with
the middle-degree kernel; other degrees vanish.
The quadratic forms of the Hamiltonians have common domain
$H^1\cap L^2(R^6\,dV)$ on forms. The variation in their coefficients
is bounded by
$C(|\epsilon|R^4+|\epsilon|^2R^2+|\epsilon|)$, which is small relative
to the common confining form as $\epsilon\to0$. Form continuity
gives norm-resolvent continuity. Compact resolvent and the constant
nine-dimensional kernel give a local spectral gap and continuous
kernel projections. A single cutoff enclosing all critical points
in this neighborhood makes the homotopy representatives in
\eqref{caps:cutoff} continuous as well.
\end{proof}

At zero, combining field rotation with target rotation gives
characters $e^{i(d-6)\theta}$ on the polynomial-degree $d$ classes.
The degrees $0,2,4$ in the determinant-character sector are distinct,
so its Gram matrix is diagonal. Combining continuity with
\eqref{caps:homogeneous}, the radial shape metric satisfies
\begin{equation}\label{end:small-rho}
g_E(\rho)=D_\rho\left[
 \operatorname{diag}(a_0,\lambda,\lambda^2/a_0)+o(1)\right]D_\rho,
\qquad
D_\rho=\operatorname{diag}(\sqrt\rho,1,\rho^{-1/2}),
\quad \rho=1/c\downarrow0.
\end{equation}
The assertion concerns the rescaled matrix. It does not assert that
every unscaled off-diagonal entry tends to zero.
For the identity class this gives
$H(c,c)\sim a_0/c$. Direct real scaling gives
$\mathcal C(c)\sim b_0/\sqrt c$, where
\[
b_0=\frac1{2\pi}\int_{\R^2}
 e^{-(Z^4+W^4+Z^2W^2)/4}\,dZ\,dW>0.
\]
Consequently $H(c,c)/\mathcal C(c)^2\to a_0/b_0^2$; no value of $a_0$ has
been inferred from the raw period.

\subsection{The massive endpoint without a uniform gap assumption}
For $c>0$, use $Z=\sqrt c\,A$, $W=\sqrt c\,B$ instead. Then
$f_c=\phi_c^*(\rho F)$ with $\rho=1/c$ and $F=f_1$.
The frame obeys
$E=\rho\,\phi_c^*((1,S,T)\Omega_Z)$, so
$g_E=\rho^2g_{\mathrm{ref}}$. All nine critical points of $F$
are now fixed and Morse.

Let $H_v$ be the complex Hessian at a critical point. A unitary
Takagi factorization gives positive singular values and local
holomorphic Morse coordinates with unit-modulus Jacobian $k_v$.
The product Gaussians of Section~\ref{sec:caps}, corrected by the
compact homotopy outside disjoint small balls, give exact-class
representatives for the primitive idempotents at twists $\pm\rho F$.
Their supports are disjoint and their squared norms satisfy
\[
N_v^\pm=\frac{4\pi^2}{\rho^2|\det H_v|}
                 (1+O(\rho^{-1/2})).
\]
The relative error follows from the first variation of the
Euclidean metric in Morse coordinates on the scale
$|\zeta|\asymp\rho^{-1/2}$; cutoff corrections are exponentially
small. Their opposite-twist pairing is exactly diagonal, with
entries
$k_{v,\rho}=4\pi^2/(\rho^2\det H_v)$.

\begin{proposition}\label{end:squeeze}
In the primitive-idempotent frame,
\[
\rho^2G_{\mathrm{ref}}
 =4\pi^2\operatorname{diag}_v(|\det H_v|^{-1})
                         +O(\rho^{-1/2}).
\]
\end{proposition}
\begin{proof}
Harmonic representatives minimize norm in their class, giving
$G_{\mathrm{ref}}\le\operatorname{diag}N_v^+$.
For a source with coefficients $z_v$, pair with the compact
opposite-twist representative of coefficients
$\overline{z_vk_{v,\rho}}/N_v^-$. Cauchy--Schwarz and exact descent
give the lower matrix bound
$\operatorname{diag}(|k_{v,\rho}|^2/N_v^-)\le G_{\mathrm{ref}}$.
The two bounds have the same leading term.
\end{proof}
This proof needs no spectral gap uniform in $\rho$.
It establishes the metric limit, not a derivative expansion or
Stokes data.

The three orbit sums of $|\det H_v|^{-1}$ are $16,64,48$.
The evaluation matrix from the polynomial frame to these orbits is
\[
\mathsf V=\begin{pmatrix}1&0&0\\1&-1&0\\1&-4/3&4/9\end{pmatrix}.
\]
Thus the other endpoint is completely fixed:
\begin{equation}\label{end:large-rho}
g_E(\rho)\longrightarrow g_\infty
=4\pi^2
\begin{pmatrix}
128&-128&64/3\\
-128&448/3&-256/9\\
64/3&-256/9&256/27
\end{pmatrix},
\qquad \rho\to\infty.
\end{equation}
In particular,
$H(c,c)=512\pi^2+O(\sqrt c)$ as $c\downarrow0$.
The exactly quadratic theory at $c=0$ has identity norm
$64\pi^2$. The difference is explained by the two escaping orbits,
whose limiting contributions are $256\pi^2$ and $192\pi^2$.
It is not legitimate to replace the small-positive-$c$ nine-vacuum
limit by the one-vacuum Gaussian model.

\subsection{Uniqueness under the physical chiral equation}
The following statement is conditional on the exact identification
of the physical connection and chiral equation in
\eqref{shape:ttstar}. Endpoint estimates alone do not establish that
identification.

\begin{theorem}\label{end:unique}
For the same Higgs matrix $C=R/16$, there is at most one positive
$C^2$ metric on the punctured shape plane satisfying
\eqref{shape:ttstar}, \eqref{end:small-rho}, and
\eqref{end:large-rho}. The constant $a_0>0$ in the first endpoint
need not be prescribed. For nonradial solutions, the endpoint
conditions are required uniformly on circles.
\end{theorem}
\begin{proof}
For two metrics put $F_{12}=\tr(g_1^{-1}g_2)$. At a point choose
a holomorphic normal frame for $g_1$, then a constant unitary
change with $g_2=\Lambda=\operatorname{diag}(\lambda_i)$.
Writing $X=\partial g_2$, direct substitution of the two metric
equations gives
\[
\bar\partial\partial F_{12}
 =\norm{\Lambda^{-1/2}X}_{\mathrm{HS}}^2
  +\sum_{i,j}\frac{(\lambda_i-\lambda_j)^2}{\lambda_j}|C_{ij}|^2
 \ge0.
\]
Therefore $d=F_{12}+F_{21}-6$ is nonnegative and subharmonic.
At the small endpoint its limit is
$2(a_2/a_1+a_1/a_2-2)$, where $a_i$ are the two possible constants,
so it is bounded. At infinity it tends to zero because both
metrics tend to $g_\infty$.
On $\epsilon<|u|<R$, bound $d$ by the harmonic function
\[
\delta_R+(K-\delta_R)
           \frac{\log(R/|u|)}{\log(R/\epsilon)},
\]
where $K$ bounds the inner boundary and
$\delta_R=\sup_{|u|=R}d$. Let $\epsilon\downarrow0$ and then
$R\to\infty$. This proves $d=0$ and hence $g_1=g_2$.
\end{proof}
The physical harmonic metric exists throughout the punctured
parameter plane by the local uniform form estimates, the
constant-dimensional kernel and fixed compact representatives.
If it satisfies the stated chiral equation in this frame, it
supplies existence as well as uniqueness. In that case the
massive endpoint selects $a_0$ implicitly. Neither its numerical
value nor a new explicit solution of the nonlinear radial equation
is asserted.
